%!TEX program = xelatex
\documentclass[11pt, a4paper]{article}

% AA LINGUA E GESTIONE FONT (Specifica per XeLaTeX)
\usepackage{fontspec}
\setmainfont{Latin Modern Roman}
\usepackage[italian]{babel}

% AA LAYOUT DI PAGINA E TIPOGRAFIA (Ottimizzato per Rilegatura)
\usepackage[left=2.4cm, right=1.6cm, top=2cm, bottom=2cm]{geometry}

% AA Spaziatura tra paragrafi controllata (compatta ed evita salti enormi)
\usepackage[skip=3pt plus 1pt minus 1pt]{parskip}
\setlength{\parindent}{0pt}

% AA Microtype per XeLaTeX
\usepackage[protrusion=true]{microtype}
\raggedbottom

% AA GESTIONE DEI COLORI
\usepackage{xcolor}

% AA MATEMATICA E SIMBOLI
\usepackage{amsmath}
\usepackage{mathtools}
\usepackage{silence}
\usepackage[warnings-off={mathtools-overbracket,mathtools-colon}]{unicode-math}
\setmathfont{Latin Modern Math}

% BB OTTIMIZZAZIONE MATRICI LUNGHE PER LA STAMPA
\setlength{\arraycolsep}{10pt}

% AA Spaziature verticali attorno alle formule matematiche
\setlength{\abovedisplayskip}{3pt plus 1pt minus 1pt}
\setlength{\belowdisplayskip}{3pt plus 1pt minus 1pt}
\setlength{\abovedisplayshortskip}{0pt plus 1pt}
\setlength{\belowdisplayshortskip}{2pt plus 1pt minus 1pt}

% AA PERSONALIZZAZIONE TITOLI
\usepackage{titlesec}
\titleformat{\section}{\normalfont\Large\bfseries\color{blue!80!black}}{\thesection}{1em}{}
\titleformat{name=\section,numberless}{\normalfont\Large\bfseries\color{blue!80!black}}{}{0pt}{}

\titleformat{\subsection}{\normalfont\large\bfseries}{\thesubsection}{1em}{}
\titleformat{name=\subsection,numberless}{\normalfont\large\bfseries}{}{0pt}{}

\titleformat{\subsubsection}{\normalfont\normalsize\bfseries}{\thesubsubsection}{1em}{}
\titleformat{name=\subsubsection,numberless}{\normalfont\normalsize\bfseries}{}{0pt}{}

\titlespacing*{\section}{0pt}{1.2ex plus 0.4ex minus 0.2ex}{0.6ex plus 0.2ex}
\titlespacing*{\subsection}{0pt}{1ex plus 0.3ex minus 0.2ex}{0.5ex plus 0.2ex}
\titlespacing*{\subsubsection}{0pt}{0.8ex plus 0.2ex minus 0.1ex}{0.4ex plus 0.1ex}

\titleformat{\paragraph}[block]{\normalfont\normalsize\bfseries\color{black!80!black}\raggedright}{}{0pt}{}
\titlespacing*{\paragraph}{0pt}{1.2ex plus 0.3ex minus 0.2ex}{0.4ex}

% AA INTEGRAZIONE CODICE MATLAB
\usepackage{matlab}

% BB Riconfigurazione Automatica Titoli Esportati da MATLAB
\renewcommand{\matlabheading}[1]{\section*{#1}}
\renewcommand{\matlabheadingtwo}[1]{\subsection*{#1}}
\renewcommand{\matlabheadingthree}[1]{\subsubsection*{#1}}

% BB Formule E Matrici Simboliche
\renewenvironment{matlabsymbolicoutput}{%
  \par\vspace{2pt}%
  \begingroup
  \centering
}{%
  \par\vspace{2pt}%
  \endgroup
}

% BB Output Testuale Console
\renewenvironment{matlaboutput}{%
  \par\vspace{1pt}%
  \Verbatim[xleftmargin=10pt, formatcom=\color{black!75}, fontsize=\footnotesize]%
}{%
  \endVerbatim
  \vspace{1pt}%
}

% BB Box Codice (Stile Moderno con Barra Laterale Colorata)
\lstset{
  backgroundcolor={},
  aboveskip=0pt,
  belowskip=0pt,
  basicstyle=\ttfamily\footnotesize,
  xleftmargin=0pt
}

\mdfdefinestyle{matlabcode}{%
  topline=false,
  bottomline=false,
  rightline=false,
  leftline=true,
  linewidth=2.5pt,
  linecolor=blue!75!black,
  backgroundcolor=black!3,
  innertopmargin=4pt,
  innerbottommargin=4pt,
  innerleftmargin=9pt,
  innerrightmargin=6pt,
  skipabove=3pt,
  skipbelow=3pt
}

\allowdisplaybreaks

% AA LINK, METADATI ED ELENCHI
\usepackage{hyperref}
\hypersetup{
    colorlinks=true,
    linkcolor=blue!80!black,
    urlcolor=blue!80!black,
    citecolor=blue!80!black,
    pdftitle={Cinematica Inversa Manipolatore Stanford},
    pdfauthor={Francesco Lo Verde},
}
\usepackage{bookmark}

% BB Spaziatura Elenchi (Arieggiata per la Stampa)
\usepackage{enumitem}
\setlist{itemsep=3.5pt plus 1pt minus 1pt, topsep=4pt plus 1pt minus 1pt, partopsep=0pt, parsep=1pt}


\author{Francesco Lo Verde}
\begin{document}

\matlabtitle{Manipolatore Stanford: Disaccoppiamento Polso - Struttura Portante}

\begin{par}
\begin{flushleft}
Il \textbf{disaccoppiamento }permette di separare il problema in parti piu semplici:
\end{flushleft}
\end{par}

\begin{itemize}

   \item{\begin{flushleft} \textbf{Cinematica Posizione: }Trovare i valori dei primi tre giunti $(q_1 ,q_2 ,q_3 )$ che posizionano il centro del polso nel punto desiderato \end{flushleft}}
   \item{\begin{flushleft} \textbf{Cinematica Orientamento: }Trovare i valori degli ultimi tre giunti $(q_4 ,q_5 ,q_6 )$ per ottenere l'orientamento desiderato dell'end effector avendo già posizionato il polso \end{flushleft}}
\end{itemize}

\begin{par}
\begin{flushleft}
Si definiscono tutte le variabili simboliche che includono le variabili di giunto $(q_1 ,\ldots,q_6 )$ e i parametri geometrici costanti del manipolatore, come le lunghezze dei link $d_2$ e $d_6$
\end{flushleft}
\end{par}

\begin{matlabcode}
clc; clear; close all;
% Variabili Simboliche Struttura
syms q1 q2 q3 d2 real;
% Variabili Simboliche Polso
syms q4 q5 q6 d6 real;
\end{matlabcode}

\begin{par}
\begin{flushleft}
Si definiscono i parametri di Denavit-Hartenberg per la struttura portante (\textbf{Giunti 1 - 3}) e per il polso (\textbf{Giunti 4 - 6}) che descrivono la relazione cinematica tra i link consecutivi.
\end{flushleft}
\end{par}

\begin{matlabcode}
% Definizione Tabelle DH Struttura
DH_Table_Strut = [q1, 0, -pi/2, 0;
                  q2, d2, pi/2, 0;
                  0, q3, 0, 0];

% Definizione Tabelle DH Polso
DH_Table_Polso = [q4, 0, -pi/2, 0;
                  q5, 0, pi/2, 0;
                  q6, d6, 0, 0];
\end{matlabcode}

\begin{par}
\begin{flushleft}
Utilizzando i parametri DH, si calcolano le \textbf{matrici di trasformazione omogenea} per ogni link tramite la funzione \texttt{Compute\_DH\_Transformations.}
\end{flushleft}
\end{par}

\begin{matlabcode}
% Calcolo Matrici Trasformazione E Rotazione Struttura
[T_Mat_Strut, R_Mat_Strut] = Compute_DH_Transformations(DH_Table_Strut);

% Calcolo Matrici Trasformazione E Rotazione Polso
[T_Mat_Polso, R_Mat_Polso] = Compute_DH_Transformations(DH_Table_Polso);
\end{matlabcode}

\begin{par}
\begin{flushleft}
A questo punto si possono comporre le matrici individuali per ottenere le trasformazioni totali:
\end{flushleft}
\end{par}

\begin{itemize}

   \item{\begin{flushleft} \texttt{T03}: Dalla base al centro del polso . \end{flushleft}}
   \item{\begin{flushleft} \texttt{T36}: dal centro del polso all'end-effector \end{flushleft}}
\end{itemize}

\begin{matlabcode}
fprintf('Matrice di trasformazione omogenea ^{%d}T_(%d):', 1, 3);
T03 = simplify(T_Mat_Strut(:,:,1) * T_Mat_Strut(:,:,2) * T_Mat_Strut(:,:,3));
disp(T03);
\end{matlabcode}
\begin{matlaboutput}
Matrice di trasformazione omogenea ^{1}T_(3):
\end{matlaboutput}
\begin{matlabsymbolicoutput}
\hskip1em $\displaystyle \left(\begin{array}{cccc}
\cos \left(q_1 \right)\,\cos \left(q_2 \right) & -\sin \left(q_1 \right) & \cos \left(q_1 \right)\,\sin \left(q_2 \right) & q_3 \,\cos \left(q_1 \right)\,\sin \left(q_2 \right)-d_2 \,\sin \left(q_1 \right)\\
\cos \left(q_2 \right)\,\sin \left(q_1 \right) & \cos \left(q_1 \right) & \sin \left(q_1 \right)\,\sin \left(q_2 \right) & d_2 \,\cos \left(q_1 \right)+q_3 \,\sin \left(q_1 \right)\,\sin \left(q_2 \right)\\
-\sin \left(q_2 \right) & 0 & \cos \left(q_2 \right) & q_3 \,\cos \left(q_2 \right)\\
0 & 0 & 0 & 1
\end{array}\right)$
\end{matlabsymbolicoutput}
\newpage
\begin{matlabcode}
fprintf('Matrice di trasformazione omogenea ^{%d}T_(%d):', 3, 6);
T36 = simplify(T_Mat_Polso(:,:,1) * T_Mat_Polso(:,:,2) * T_Mat_Polso(:,:,3));
disp(T36);
\end{matlabcode}
\begin{matlaboutput}
Matrice di trasformazione omogenea ^{3}T_(6):
\end{matlaboutput}
\resizebox{\linewidth}{!}{%
$
\begin{pmatrix}
    c\left(q_4 \right)\,c\left(q_5 \right)\,c\left(q_6 \right)-s\left(q_4 \right)\,s\left(q_6 \right) & -c\left(q_6 \right)\,s\left(q_4 \right)-c\left(q_4 \right)\,c\left(q_5 \right)\,s\left(q_6 \right) & c\left(q_4 \right)\,s\left(q_5 \right) & d_6 \,c\left(q_4 \right)\,s\left(q_5 \right)\\
    c\left(q_4 \right)\,s\left(q_6 \right)+c\left(q_5 \right)\,c\left(q_6 \right)\,s\left(q_4 \right) & c\left(q_4 \right)\,c\left(q_6 \right)-c\left(q_5 \right)\,s\left(q_4 \right)\,s\left(q_6 \right) & s\left(q_4 \right)\,s\left(q_5 \right) & d_6 \,s\left(q_4 \right)\,s\left(q_5 \right)\\
    -c\left(q_6 \right)\,s\left(q_5 \right) & s\left(q_5 \right)\,s\left(q_6 \right) & c\left(q_5 \right) & d_6 \,c\left(q_5 \right)\\
    0 & 0 & 0 & 1
\end{pmatrix}
$}
\begin{par}
\begin{flushleft}
La terna 3 del \textbf{manipolatore sferico} coincide con la terna 3 del \textbf{polso sferico}, la funzione cinematica diretta può essere ottenuta semplicemente per composizione delle matrici di trasformazione come${\,}^1 T_6 ={\,}^1 T_3 \cdot {\,}^3 T_6$
\end{flushleft}
\end{par}

\begin{matlabcode}
T06 = simplify(T03 * T36);
% fprintf('Matrice di trasformazione omogenea ^{%d}T_(%d):', 0, 6);
% disp(T16);
\end{matlabcode}
\begin{par}
\begin{flushleft}
Il vettore $d_{36}$ rappresenta il \textbf{vettore di traslazione} che descrive la posizione dell'origine del sistema di riferimento dell'end-effector rispetto all'origine del sistema di riferimento del polso.\\
In pratica, è il \underline{vettore di posizione dell'end-effector visto dal centro del polso}.
\end{flushleft}
\end{par}

\begin{matlabcode}
d36 = T36(1:3, 4);
fprintf('Vettore d_{36}');
disp(d36);
\end{matlabcode}
\begin{matlaboutput}
Vettore d_{36}
\end{matlaboutput}
\begin{matlabsymbolicoutput}
\hskip1em $\displaystyle \left(\begin{array}{c}
d_6 \,\cos \left(q_4 \right)\,\sin \left(q_5 \right)\\
d_6 \,\sin \left(q_4 \right)\,\sin \left(q_5 \right)\\
d_6 \,\cos \left(q_5 \right)
\end{array}\right)$
\end{matlabsymbolicoutput}


\vspace{1em}

\matlabheading{Cinematica Inversa Posizione}

\begin{par}
\begin{flushleft}
Si deve verificare che il vettore che collega il centro del polso all'end-effector, quando espresso nelle coordinate dell'end-effector (terna 6), sia costante. Questo vettore \texttt{d1} si calcola come:
\end{flushleft}
\end{par}

\begin{par}
$\exists \;d_0 ,\,d_1 \in {\mathbb{R}}^3 \;\mid \;d_{36} =^3 R_6 d_1 +d_0$
\end{par}

\begin{par}
\begin{flushleft}
Dato che $d_{36}$ non presenta termini indipendenti dalle variabili di giunto, questo ci dice che non c'è un offset fisso dal polso, e quindi non c'è uno spostamento strutturale fisso, visto dal polso e di conseguenza allora $d_0 =0$
\end{flushleft}
\end{par}

\[
    d_{36} = {}^{3}R_{6} \; d_{1} \qquad \longrightarrow \qquad d_{1} =  =\bigl({}^{3}R_{6}\bigr)^{T} d_{36}
\]

\begin{par}
\begin{flushleft}
Poichè $^3 R_6$ è una matrice di rotazione, si ha che ${\bigl(^3 R_6 \bigr)}^{-1} ={\bigl(^3 R_6 \bigr)}^T$
\end{flushleft}
\end{par}

\begin{matlabcode}
R36 = simplify(R_Mat_Polso(:,:,1) * R_Mat_Polso(:,:,2) * R_Mat_Polso(:,:,3));
R36_T = R36';
d1 = simplify(R36_T * d36);
fprintf('Vettore d_{1}');
disp(d1);
\end{matlabcode}
\begin{matlaboutput}
Vettore d_{1}
\end{matlaboutput}
\begin{matlabsymbolicoutput}
\hskip1em $\displaystyle \left(\begin{array}{c}
0\\
0\\
d_6
\end{array}\right)$
\end{matlabsymbolicoutput}
\newpage
\begin{par}
\begin{flushleft}
Questo è il vettore che collega il centro del polso (origine di $RF_3$) all'end-effector (origine di $RF_6$), espresso in coordinate del sistema $RF_6$.
\end{flushleft}
\end{par}

\begin{par}
\begin{flushleft}
Sapendo che la posizione desiderata dell'end-effector è $d_{06}$, si può scrivere
\end{flushleft}
\end{par}

\begin{par}
$d_{06} =d_{03} +^0 R_6 \;d_1$
\end{par}

\begin{par}
\begin{flushleft}
Cioè, il vettore che va dalla \textbf{base all'end effector} è uguale al vettore che va dalla \textbf{base al polso} più il vettore che va dal \textbf{polso all'end effector. }Quindi, dalla precedente si può dire che la posizione del centro del polso $d_{03}$ può essere calcolata come $d_{03} =d_{06} -^0 R_6 \;d_1$, dove $d_{06}$ è il vettore obiettivo cioè la posizione finale che deve avere l'end effector
\end{flushleft}
\end{par}

\begin{matlabcode}
syms dx dy dz real;

% Posa Desiderata (INPUT DEL PROBLEMA)
d06 = [dx; dy; dz];

% Orientamento Desiderato Tramite Matrice Di Rotazione (INPUT DEL PROBLEMA)
R06_End = sym('R', [3, 3], 'real');

fprintf('Posizione Desiderata (d06):');
disp(d06);
\end{matlabcode}
\begin{matlaboutput}
Posizione Desiderata (d06):
\end{matlaboutput}
\begin{matlabsymbolicoutput}
\hskip1em $\displaystyle \left(\begin{array}{c}
\textrm{dx}\\
\textrm{dy}\\
\textrm{dz}
\end{array}\right)$
\end{matlabsymbolicoutput}
\begin{matlabcode}
fprintf('Orientamento Desiderato (R06_End):');
disp(R06_End);
\end{matlabcode}
\begin{matlaboutput}
Orientamento Desiderato (R06_End):
\end{matlaboutput}
\begin{matlabsymbolicoutput}
\hskip1em $\displaystyle \left(\begin{array}{ccc}
R_{1,1}  & R_{1,2}  & R_{1,3} \\
R_{2,1}  & R_{2,2}  & R_{2,3} \\
R_{3,1}  & R_{3,2}  & R_{3,3}
\end{array}\right)$
\end{matlabsymbolicoutput}
\begin{matlabcode}

d03 = T03(1:3, 4);
fprintf('Posizione Calcolata Struttura Portante (d03):');
disp(d03);
\end{matlabcode}
\begin{matlaboutput}
Posizione Calcolata Struttura Portante (d03):
\end{matlaboutput}
\begin{matlabsymbolicoutput}
\hskip1em $\displaystyle \left(\begin{array}{c}
q_3 \,\cos \left(q_1 \right)\,\sin \left(q_2 \right)-d_2 \,\sin \left(q_1 \right)\\
d_2 \,\cos \left(q_1 \right)+q_3 \,\sin \left(q_1 \right)\,\sin \left(q_2 \right)\\
q_3 \,\cos \left(q_2 \right)
\end{array}\right)$
\end{matlabsymbolicoutput}

\newpage
\matlabheadingtwo{Valore $q_3$}

\begin{par}
\begin{flushleft}
Si indica con $d_{hat}$ il vettore che rappresenta la posizione desiderata del centro del polso rispetto al sistema di riferimento di base. Questo si può calcolare come:
\end{flushleft}
\end{par}

\begin{par}
$\hat{d} =d_{06} -^0 R_6 \cdot d_1$
\end{par}

\begin{par}
\begin{flushleft}
Questo vettore d\_\{hat\} è considerato \textbf{noto}, in quanto $d_{06}$ (la posizione finale) e $^0 R_6$ (l'orientamento finale) sono \textbf{gli input problema} di cinematica inversa.
\end{flushleft}
\end{par}

\begin{matlabcode}
% Box di Calcolo d_hat
% d_hat = d06_des - R06_des * d1_vec

% Calcolo posizione target del polso
d_hat = d06 - R06_End * d1;
fprintf('Posizione Desiderata Struttura Portante (d_hat):');
disp(d_hat);
\end{matlabcode}
\begin{matlaboutput}
Posizione Desiderata Struttura Portante (d_hat):
\end{matlaboutput}
\begin{matlabsymbolicoutput}
\hskip1em $\displaystyle \left(\begin{array}{c}
\textrm{dx}-R_{1,3} \,d_6 \\
\textrm{dy}-R_{2,3} \,d_6 \\
\textrm{dz}-R_{3,3} \,d_6
\end{array}\right)$
\end{matlabsymbolicoutput}
\begin{matlabcode}
% Avendo come input un vettore e un orientamento numerici, questo sara un vettore numerico
\end{matlabcode}

\begin{par}
\begin{flushleft}
A questo punto si hanno due rappresentazioni della posizione del centro del polso:
\end{flushleft}
\end{par}

\begin{itemize}

   \item{\begin{flushleft} $\hat{d}$: che è un vettore noto in quanto calcolato tramite gli \textbf{input del problema} \end{flushleft}}
   \item{\begin{flushleft} $d_{03}$: che è il vettore di posizione calcolato dalla cinematica diretta e dipende da $q1,q2,q3$ \end{flushleft}}
\end{itemize}

\begin{par}
\begin{flushleft}
L'end-effector deve poter raggiungere la posa, quindi una condizione necessaria è che la lunghezza (o norma) del vettore target $\hat{d}$ sia uguale alla lunghezza del vettore $d_{03}$. Quindi, uguagliando le loro norme al quadrato, si  ottiene:
\end{flushleft}
\end{par}

\begin{matlabcode}
norm_sq_d_hat = norm(d_hat)^2;
norm_sq_d03 = norm(d03)^2;

norm_sq_d_hat = simplify(norm_sq_d_hat);
norm_sq_d03 = simplify(norm_sq_d03);

fprintf('Norma del vettore target (d_hat):');

disp(norm_sq_d_hat);
\end{matlabcode}
\begin{matlaboutput}
Norma del vettore target (d_hat):
\end{matlaboutput}
\begin{matlabsymbolicoutput}
\hskip1em $\displaystyle {{\left(\textrm{dx}-R_{1,3} \,d_6 \right)}}^2 +{{\left(\textrm{dy}-R_{2,3} \,d_6 \right)}}^2 +{{\left(\textrm{dz}-R_{3,3} \,d_6 \right)}}^2 $
\end{matlabsymbolicoutput}

\begin{par}
\begin{flushleft}
Come detto prima $\hat{d}$  è un vettore noto, la sua norma al quadrato sarà un valore numerico. Poichè si stanno eseguendo i calcoli in maniera simbolica, per evitare risultati simbolici troppo lunghi e complessi chiamo questo valore \texttt{dHatNormSq}
\end{flushleft}
\end{par}

\begin{matlabcode}
syms dHatNormSq real;

fprintf('Norma al quadrato del vettore calcolato (d_03):');

disp(norm_sq_d03);
\end{matlabcode}
\begin{matlaboutput}
Norma al quadrato del vettore calcolato (d_03):
\end{matlaboutput}
\begin{matlabsymbolicoutput}
\hskip1em $\displaystyle {d_2 }^2 +{q_3 }^2 $
\end{matlabsymbolicoutput}
\newpage
\begin{matlabcode}
eq_q3 = (dHatNormSq == norm_sq_d03);

fprintf('Equazione per trovare q3:');
disp(eq_q3);
\end{matlabcode}
\begin{matlaboutput}
Equazione per trovare q3:
\end{matlaboutput}
\begin{matlabsymbolicoutput}
\hskip1em $\displaystyle \textrm{dHatNormSq}={d_2 }^2 +{q_3 }^2 $
\end{matlabsymbolicoutput}
\begin{matlabcode}

% Risolviamo per q3
sol_q3 = solve(eq_q3, q3, ReturnConditions=true);
q3_pos = sol_q3.q3(1);
q3_neg = sol_q3.q3(2);

fprintf('Soluzioni per q3:');
fprintf('q3_1 = ');
disp(q3_pos);
\end{matlabcode}
\begin{matlaboutput}
Soluzioni per q3:
q3_1 =
\end{matlaboutput}
\begin{matlabsymbolicoutput}
\hskip1em $\displaystyle \sqrt{\textrm{dHatNormSq}-{d_2 }^2 }$
\end{matlabsymbolicoutput}
\begin{matlabcode}
fprintf('q3_2 =');
disp(q3_neg);
\end{matlabcode}
\begin{matlaboutput}
q3_2 =
\end{matlaboutput}
\begin{matlabsymbolicoutput}
\hskip1em $\displaystyle -\sqrt{\textrm{dHatNormSq}-{d_2 }^2 }$
\end{matlabsymbolicoutput}
\begin{matlabcode}
fprintf('Entrambe le soluzioni sono valide se:');
disp(sol_q3.conditions(2));
\end{matlabcode}
\begin{matlaboutput}
Entrambe le soluzioni sono valide se:
\end{matlaboutput}
\begin{matlabsymbolicoutput}
\hskip1em $\displaystyle {d_2 }^2 \le \textrm{dHatNormSq}$
\end{matlabsymbolicoutput}

\begin{par}
\begin{flushleft}
Queste due soluzioni corrispondono a due diverse configurazioni del braccio:
\end{flushleft}
\end{par}

\begin{itemize}
   \item{\begin{flushleft} \textbf{Gomito Alto} (Soluzione Positiva): Il gomito è sopra la linea che collega la base all'end-effector \end{flushleft}}
   \item{\begin{flushleft} \textbf{Gomito Basso} (Soluzione Negativa): Il gomito è sotto la linea che collega la base all'end-effector \end{flushleft}}
\end{itemize}

\begin{par}
\begin{flushleft}
Ugualiando $\hat{d}$ (\textbf{Vettore posizione target polso}) e $d_{03}$ (\textbf{Vettore posizione polso cinematica diretta}) si ottiene il sistema
\end{flushleft}
\end{par}

\begin{matlabcode}
fprintf('Sistema da risolvere per trovare q1 e q2');
eq_system = d03 == d_hat;
disp(eq_system);
\end{matlabcode}
\begin{matlaboutput}
Sistema da risolvere per trovare q1 e q2
\end{matlaboutput}
\begin{matlabsymbolicoutput}
\hskip1em $\displaystyle \left(\begin{array}{c}
q_3 \,\cos \left(q_1 \right)\,\sin \left(q_2 \right)-d_2 \,\sin \left(q_1 \right)=\textrm{dx}-R_{1,3} \,d_6 \\
d_2 \,\cos \left(q_1 \right)+q_3 \,\sin \left(q_1 \right)\,\sin \left(q_2 \right)=\textrm{dy}-R_{2,3} \,d_6 \\
q_3 \,\cos \left(q_2 \right)=\textrm{dz}-R_{3,3} \,d_6
\end{array}\right)$
\end{matlabsymbolicoutput}
\newpage
\matlabheadingtwo{Valore $q_2$}

\begin{par}
\begin{flushleft}
Si nota che, dalla terza equazione si può trovare $\cos (q_2 )$
\end{flushleft}
\end{par}

\begin{matlabcode}
syms c2 real;
eq_cosq2 = eq_system(3);
fprintf('Terza Equazione:');
disp(eq_cosq2);
\end{matlabcode}
\begin{matlaboutput}
Terza Equazione:
\end{matlaboutput}
\begin{matlabsymbolicoutput}
\hskip1em $\displaystyle q_3 \,\cos \left(q_2 \right)=\textrm{dz}-R_{3,3} \,d_6 $
\end{matlabsymbolicoutput}
\begin{matlabcode}

cosq2_eq = isolate(eq_cosq2, cos(q2));
fprintf('Equazione cos(q2):');
disp(cosq2_eq);
\end{matlabcode}
\begin{matlaboutput}
Equazione cos(q2):
\end{matlaboutput}
\begin{matlabsymbolicoutput}
\hskip1em $\displaystyle \cos \left(q_2 \right)=\frac{\textrm{dz}-R_{3,3} \,d_6 }{q_3 }$
\end{matlabsymbolicoutput}

\begin{par}
\begin{flushleft}
Poiché $\sin^2 (q_2 )+\cos^2 (q_2 )=1$, si hanno due possibili soluzioni per $\sin (q_2 )$:
\end{flushleft}
\end{par}

\begin{par}
$\sin (q_2 )=\pm \sqrt{1-\cos (q_2 )^2 \;\;}$
\end{par}

\begin{matlabcode}
cosq2 = rhs(cosq2_eq);

sinq2_pos = sqrt(1 - cosq2^2);
sinq2_neg = -sqrt(1 - cosq2^2);

q2_neg = atan2(sinq2_neg, cosq2);
q2_pos = atan2(sinq2_pos, cosq2);

fprintf('Valore q2 Per Valore Seno Negativo:');

disp(q2_neg);
\end{matlabcode}
\begin{matlaboutput}
Valore q2 Per Valore Seno Negativo:
\end{matlaboutput}

\begin{matlabsymbolicoutput}
\hskip1em $\displaystyle \textrm{atan2}\left(-\sqrt{1-\frac{{{\left(\textrm{dz}-R_{3,3} \,d_6 \right)}}^2 }{{q_3 }^2 }},\frac{\textrm{dz}-R_{3,3} \,d_6 }{q_3 }\right)$
\end{matlabsymbolicoutput}
\begin{matlabcode}
fprintf('Valore q2 Per Valore Seno Positivo:');
disp(q2_pos);
\end{matlabcode}
\begin{matlaboutput}
Valore q2 Per Valore Seno Positivo:
\end{matlaboutput}

\begin{matlabsymbolicoutput}
\hskip1em $\displaystyle \textrm{atan2}\left(\sqrt{1-\frac{{{\left(\textrm{dz}-R_{3,3} \,d_6 \right)}}^2 }{{q_3 }^2 }},\frac{\textrm{dz}-R_{3,3} \,d_6 }{q_3 }\right)$
\end{matlabsymbolicoutput}

\newpage
\matlabheadingtwo{Valore $q_1$}

\begin{par}
\begin{flushleft}
Si nota che, la prima e la seconda equazione del precedente sistema formano un sistema utile per trovare$\cos (q_1 )$
\end{flushleft}
\end{par}

\begin{matlabcode}
syms c1 s1 real;
eq_q1 = eq_system(1:2);
disp(eq_q1);
\end{matlabcode}
\begin{matlabsymbolicoutput}
\hskip1em $\displaystyle \left(\begin{array}{c}
q_3 \,\cos \left(q_1 \right)\,\sin \left(q_2 \right)-d_2 \,\sin \left(q_1 \right)=\textrm{dx}-R_{1,3} \,d_6 \\
d_2 \,\cos \left(q_1 \right)+q_3 \,\sin \left(q_1 \right)\,\sin \left(q_2 \right)=\textrm{dy}-R_{2,3} \,d_6
\end{array}\right)$
\end{matlabsymbolicoutput}
\begin{matlabcode}

eq_q1_sym = subs(eq_q1, [cos(q1), sin(q1)], [c1, s1]);

sol_q1 = solve(eq_q1_sym, [c1, s1]);
sinq1 = simplify(sol_q1.s1);
cosq1 = simplify(sol_q1.c1);

fprintf('Soluzioni Del Sistema:');

disp(sin(q1) == sinq1);

disp(cos(q1) == cosq1);
\end{matlabcode}
\begin{matlaboutput}
Soluzioni Del Sistema:
\end{matlaboutput}
\begin{matlabsymbolicoutput}
\hskip1em $\displaystyle \sin \left(q_1 \right)=-\frac{d_2 \,\textrm{dx}-\textrm{dy}\,q_3 \,\sin \left(q_2 \right)-R_{1,3} \,d_2 \,d_6 +R_{2,3} \,d_6 \,q_3 \,\sin \left(q_2 \right)}{{d_2 }^2 +{q_3 }^2 \,{\sin \left(q_2 \right)}^2 }$
\end{matlabsymbolicoutput}
\begin{matlabsymbolicoutput}
\hskip1em $\displaystyle \cos \left(q_1 \right)=\frac{d_2 \,\textrm{dy}+\textrm{dx}\,q_3 \,\sin \left(q_2 \right)-R_{2,3} \,d_2 \,d_6 -R_{1,3} \,d_6 \,q_3 \,\sin \left(q_2 \right)}{{d_2 }^2 +{q_3 }^2 \,{\sin \left(q_2 \right)}^2 }$
\end{matlabsymbolicoutput}

\begin{par}
\begin{flushleft}
A questo punto, usando la funzione arcotangente a quattro quadranti si può calcolare il valore di $q_1$
\end{flushleft}
\end{par}

\begin{matlabcode}
q1_val = atan2(sinq1, cosq1);

fprintf('Valore q1:');
disp(q1_val);
\end{matlabcode}
\begin{matlaboutput}
Valore q1:
\end{matlaboutput}
\begin{matlabsymbolicoutput}
\resizebox{\linewidth}{!}{%
$
    \textrm{atan2}\left(-\frac{d_2 \,\textrm{dx}-\textrm{dy}\,q_3 \,\sin \left(q_2 \right)-R_{1,3} \,d_2 \,d_6 +R_{2,3} \,d_6 \,q_3 \,\sin \left(q_2 \right)}{{d_2 }^2 +{q_3 }^2 \,{\sin \left(q_2 \right)}^2 },\frac{d_2 \,\textrm{dy}+\textrm{dx}\,q_3 \,\sin \left(q_2 \right)-R_{2,3} \,d_2 \,d_6 -R_{1,3} \,d_6 \,q_3 \,\sin \left(q_2 \right)}{{d_2 }^2 +{q_3 }^2 \,{\sin \left(q_2 \right)}^2 }\right)
$}
\end{matlabsymbolicoutput}


\vspace{1em}

\matlabheadingtwo{Cinematica Inversa Orientamento}

\begin{par}
\begin{flushleft}
Una volta determinati $q_1 ,q_2 ,q_3$, la matrice di rotazione $^0 R_3$ è nota, quindi si può o calcolare la matrice di rotazione del polso $^3 R_6$richiesta per raggiungere l'orientamento finale $^0 R_6$ dato che:
\end{flushleft}
\end{par}

\begin{par}
$^0 R_6 =^0 R_3 \cdot ^3 R_6$
\end{par}

\begin{par}
\begin{flushleft}
Sapendo che le matrici di rotazioni sono invertibili e l'inversa di una matrice di rotazione è uguale alla trasposta si ha che la matrice "target" per ottenere l'orientamento desiderato è
\end{flushleft}
\end{par}

\begin{par}
$^3 R_6 ={\bigl(^0 R_3 \bigr)}^{-1} \;^0 R_6 ={\bigl(^0 R_3 \bigr)}^T \;^0 R_6$
\end{par}
\newpage
\begin{par}
\begin{flushleft}
Dove
\end{flushleft}
\end{par}

\begin{itemize}

   \item{\begin{flushleft} $^0 R_6$: Rappresenta la matrice di rotazione desiderata per l'end-effector, che è un dato di input del problema di cinematica inversa. \end{flushleft}}
   \item{\begin{flushleft} ${\bigl(^0 R_3 \bigr)}^T$ Dipende da $q_1 ,q_2 ,q_3$ che sono appena stati calcolati. Sostituendo i loro valori numerici nella matrice di rotazione della struttura portante (estratta dalla matrice $^0 T_3$ ), otteniamo una matrice numerica conosciuta. \end{flushleft}}
   \item{\begin{flushleft} $^3 R_6$: Dipende dagli angoli $q_2 ,q_3 ,q_4$  che si devono trovare. \end{flushleft}}
\end{itemize}

\begin{par}
\begin{flushleft}
\textbf{Poichè si stanno facendo dei calcoli simbolici, indichiamo con }\texttt{\textbf{R03\_K}}\textbf{ (Known) la matrice che contiene i valori numerici calcolati (}\textit{\textbf{termini noti}}\textbf{) e con }\texttt{\textbf{R36\_T}}\textbf{ (Target) la matrice di rotazione estratta dalla matrice di trasformazione omogenea }$^3 T_6$\textbf{ ottenuta tramite le tabelle di Denavit-Hartenberg che dipende da }$q_4 ,q_5 ,q_6$\textbf{ (}\textit{\textbf{incognite}}\textbf{)}
\end{flushleft}
\end{par}

\begin{matlabcode}
% Valori noti a questo punto "_K"
% Valori da calcolare "_T"
% Orientamento Finale R06_K = R03_K * R36_T

% Matrice Rotazione R06 Nota (Input Problema)
R06_K = sym('R06K', [3, 3], 'real');
fprintf('Matrice Rotazione R06 Nota:');

disp(R06_K);
\end{matlabcode}
\begin{matlaboutput}
Matrice Rotazione R06 Nota:
\end{matlaboutput}
\begin{matlabsymbolicoutput}
\hskip1em $\displaystyle \left(\begin{array}{ccc}
{\textrm{R06K}}_{1,1}  & {\textrm{R06K}}_{1,2}  & {\textrm{R06K}}_{1,3} \\
{\textrm{R06K}}_{2,1}  & {\textrm{R06K}}_{2,2}  & {\textrm{R06K}}_{2,3} \\
{\textrm{R06K}}_{3,1}  & {\textrm{R06K}}_{3,2}  & {\textrm{R06K}}_{3,3}
\end{array}\right)$
\end{matlabsymbolicoutput}
\begin{matlabcode}

% Matrice Rotazione R03 Nota (Calcolata Con q1, q2, q3)
R03_K = sym('R03K', [3, 3], 'real');
fprintf('Matrice Rotazione R03 Nota:');

disp(R03_K);
\end{matlabcode}
\begin{matlaboutput}
Matrice Rotazione R03 Nota:
\end{matlaboutput}
\begin{matlabsymbolicoutput}
\hskip1em $\displaystyle \left(\begin{array}{ccc}
{\textrm{R03K}}_{1,1}  & {\textrm{R03K}}_{1,2}  & {\textrm{R03K}}_{1,3} \\
{\textrm{R03K}}_{2,1}  & {\textrm{R03K}}_{2,2}  & {\textrm{R03K}}_{2,3} \\
{\textrm{R03K}}_{3,1}  & {\textrm{R03K}}_{3,2}  & {\textrm{R03K}}_{3,3}
\end{array}\right)$
\end{matlabsymbolicoutput}
\begin{matlabcode}

% Matrice Rotazione R36 Con Incognite
R36_T = simplify(R_Mat_Polso(:,:,1) * R_Mat_Polso(:,:,2) * R_Mat_Polso(:,:,3));
fprintf('Matrice Rotazione R36 Incognite:');
disp(R36_T);
\end{matlabcode}
\begin{matlaboutput}
Matrice Rotazione R36 Incognite:
\end{matlaboutput}

\begin{matlabsymbolicoutput}
\resizebox{\linewidth}{!}{%
$
    \begin{pmatrix}
    c\left(q_4 \right)\,c\left(q_5 \right)\,c\left(q_6 \right)-s\left(q_4 \right)\,s\left(q_6 \right) & -c\left(q_6 \right)\,s\left(q_4 \right)-c\left(q_4 \right)\,c\left(q_5 \right)\,s\left(q_6 \right) & c\left(q_4 \right)\,s\left(q_5 \right)\\
    -c\left(q_6 \right)\,s\left(q_5 \right) & s\left(q_5 \right)\,s\left(q_6 \right) & c\left(q_5 \right)
    \end{pmatrix}
$}
\end{matlabsymbolicoutput}

\begin{par}
\begin{flushleft}
Con la formula $^3 R_6 ={\bigl(^0 R_3 \bigr)}^T \;^0 R_6$ si può calcolare la matrice $^3 R_6$ nota (\texttt{R36\_K}) e poi confrontandola membro a membro con $^3 R_6$ incognita si ottiene un sistema con incognite ($q_4 ,q_5 ,q_6$)
\end{flushleft}
\end{par}
\newpage
\begin{matlabcode}
R03_K_Trasp = R03_K';

R36_Calc = R03_K_Trasp * R06_K;
fprintf('Matrice Rotazione R36 Nota Calcolata:');
disp(R36_Calc);
\end{matlabcode}
\begin{matlaboutput}
Matrice Rotazione R36 Nota Calcolata:
\end{matlaboutput}
\begin{matlabsymbolicoutput}
\resizebox{\linewidth}{!}{%
$
    \begin{pmatrix}
        {\textrm{R03K}}_{1,1} \,{\textrm{R06K}}_{1,1} + {\textrm{R03K}}_{2,1} \,{\textrm{R06K}}_{2,1} +{\textrm{R03K}}_{3,1} \,{\textrm{R06K}}_{3,1}  & {\textrm{R03K}}_{1,1} \,{\textrm{R06K}}_{1,2} +{\textrm{R03K}}_{2,1} \,{\textrm{R06K}}_{2,2} +{\textrm{R03K}}_{3,1} \,{\textrm{R06K}}_{3,2}  & {\textrm{R03K}}_{1,1} \,{\textrm{R06K}}_{1,3} +{\textrm{R03K}}_{2,1} \,{\textrm{R06K}}_{2,3} +{\textrm{R03K}}_{3,1} \,{\textrm{R06K}}_{3,3} \\
        {\textrm{R03K}}_{1,2} \,{\textrm{R06K}}_{1,1} +{\textrm{R03K}}_{2,2} \,{\textrm{R06K}}_{2,1} +{\textrm{R03K}}_{3,2} \,{\textrm{R06K}}_{3,1}  & {\textrm{R03K}}_{1,2} \,{\textrm{R06K}}_{1,2} +{\textrm{R03K}}_{2,2} \,{\textrm{R06K}}_{2,2} +{\textrm{R03K}}_{3,2} \,{\textrm{R06K}}_{3,2}  & {\textrm{R03K}}_{1,2} \,{\textrm{R06K}}_{1,3} +{\textrm{R03K}}_{2,2} \,{\textrm{R06K}}_{2,3} +{\textrm{R03K}}_{3,2} \,{\textrm{R06K}}_{3,3} \\
        {\textrm{R03K}}_{1,3} \,{\textrm{R06K}}_{1,1} +{\textrm{R03K}}_{2,3} \,{\textrm{R06K}}_{2,1} +{\textrm{R03K}}_{3,3} \,{\textrm{R06K}}_{3,1}  & {\textrm{R03K}}_{1,3} \,{\textrm{R06K}}_{1,2} +{\textrm{R03K}}_{2,3} \,{\textrm{R06K}}_{2,2} +{\textrm{R03K}}_{3,3} \,{\textrm{R06K}}_{3,2}  & {\textrm{R03K}}_{1,3} \,{\textrm{R06K}}_{1,3} +{\textrm{R03K}}_{2,3} \,{\textrm{R06K}}_{2,3} +{\textrm{R03K}}_{3,3} \,{\textrm{R06K}}_{3,3}
    \end{pmatrix}
$}
\end{matlabsymbolicoutput}

\begin{par}
\begin{flushleft}
\textbf{Per semplificare la leggibilità delle formule, indichiamo con A}\texttt{\textbf{ }}\textbf{la matrice precedente e ogni suo elemento sarà }$\mathbf{A_{ij} }$
\end{flushleft}
\end{par}

\begin{matlabcode}
% Matrice R36 Calcolata Che Rinomino In A
R36_K = sym('A', [3, 3], 'real');
fprintf('Matrice Rotazione R03 Calcolata Con Termini Rinominati Con A_{ij}:');
disp(R36_K);
\end{matlabcode}
\begin{matlaboutput}
Matrice Rotazione R03 Calcolata Con Termini Rinominati Con A_{ij}:
\end{matlaboutput}
\begin{matlabsymbolicoutput}
\hskip1em $\displaystyle \left(\begin{array}{ccc}
A_{1,1}  & A_{1,2}  & A_{1,3} \\
A_{2,1}  & A_{2,2}  & A_{2,3} \\
A_{3,1}  & A_{3,2}  & A_{3,3}
\end{array}\right)$
\end{matlabsymbolicoutput}

\begin{par}
\begin{flushleft}
Quindi si ottiene:
\end{flushleft}
\end{par}

\begin{matlabcode}
eqSystem = R36_T == R36_K;
fprintf('Sistema Da Risolvere Per Trovare q4, q5, q6');
disp(eqSystem);
\end{matlabcode}
\begin{matlaboutput}
Sistema Da Risolvere Per Trovare q4, q5, q6
\end{matlaboutput}
\begin{matlabsymbolicoutput}
\resizebox{\linewidth}{!}{%
$
    \begin{pmatrix}
        c\left(q_4 \right)\,c\left(q_5 \right)\,c\left(q_6 \right)-s\left(q_4 \right)\,s\left(q_6 \right)=A_{1,1}  & -c\left(q_6 \right)\,s\left(q_4 \right)-c\left(q_4 \right)\,c\left(q_5 \right)\,s\left(q_6 \right)=A_{1,2}  & c\left(q_4 \right)\,s\left(q_5 \right)=A_{1,3} \\
        c\left(q_4 \right)\,s\left(q_6 \right)+c\left(q_5 \right)\,c\left(q_6 \right)\,s\left(q_4 \right)=A_{2,1}  & c\left(q_4 \right)\,c\left(q_6 \right)-c\left(q_5 \right)\,s\left(q_4 \right)\,s\left(q_6 \right)=A_{2,2}  & s\left(q_4 \right)\,s\left(q_5 \right)=A_{2,3} \\
        -c\left(q_6 \right)\,s\left(q_5 \right)=A_{3,1}  & s\left(q_5 \right)\,s\left(q_6 \right)=A_{3,2}  & c\left(q_5 \right)=A_{3,3}
    \end{pmatrix}
$}
\end{matlabsymbolicoutput}

\begin{par}
\begin{flushleft}
Si nota che il valore di $\cos (q_5 )$ è noto in quanto è l'elemento $(3,3)$ della matrice $^3 R_6$ e di conseguenza si può calcolare anche $\sin (q_5 )$ e quindi il valore dell'angolo di giunto
\end{flushleft}
\end{par}

\begin{matlabcode}
cosq5 = rhs(eqSystem(3,3));

sin_q5_pos = sqrt(1 - cosq5^2);
sin_q5_neg = -sqrt(1 - cosq5^2);

q5_neg = atan2(sin_q5_neg, cosq5);
q5_pos = atan2(sin_q5_pos, cosq5);

fprintf('Valore q5 Per Valore Seno Negativo:');
disp(q5_neg);
\end{matlabcode}
\begin{matlaboutput}
Valore q5 Per Valore Seno Negativo:
\end{matlaboutput}
\begin{matlabsymbolicoutput}
\hskip1em $\displaystyle \textrm{atan2}\left(-\sqrt{1-{A_{3,3} }^2 },A_{3,3} \right)$
\end{matlabsymbolicoutput}
\begin{matlabcode}
fprintf('Valore q5 Per Valore Seno Positivo:');
disp(q5_pos);
\end{matlabcode}
\begin{matlaboutput}
Valore q5 Per Valore Seno Positivo:
\end{matlaboutput}
\begin{matlabsymbolicoutput}
\hskip1em $\displaystyle \textrm{atan2}\left(\sqrt{1-{A_{3,3} }^2 },A_{3,3} \right)$
\end{matlabsymbolicoutput}
\newpage
\begin{par}
\begin{flushleft}
\textbf{Assumendo }il caso non singolare, cioè\textbf{ }$\mathbf{\sin (q_5 )\not= 0}$
\end{flushleft}
\end{par}

\begin{itemize}

   \item{\begin{flushleft} Si può definire un sistema con gli elementi $(1,3)$ e $(2,3)$ della matrice $^3 R_6$ per calcolare $q_4$ \end{flushleft}}
\end{itemize}

\begin{matlabcode}
fprintf('Sistema Da Risolvere Per Trovare q4:');
eq_cosq4 = eqSystem(1,3);
eq_sinq4 = eqSystem(2,3);
disp(eq_cosq4);
disp(eq_sinq4);
\end{matlabcode}
\begin{matlaboutput}
Sistema Da Risolvere Per Trovare q4:
\end{matlaboutput}
\begin{matlabsymbolicoutput}
\hskip1em $\displaystyle \cos \left(q_4 \right)\,\sin \left(q_5 \right)=A_{1,3} $
\end{matlabsymbolicoutput}
\begin{matlabsymbolicoutput}
\hskip1em $\displaystyle \sin \left(q_4 \right)\,\sin \left(q_5 \right)=A_{2,3} $
\end{matlabsymbolicoutput}
\begin{matlabcode}
fprintf('Soluzioni:');
cosq6_eq = isolate(eq_cosq4, cos(q4));
disp(cosq6_eq);
sinq4_eq = isolate(eq_sinq4, sin(q4));
disp(sinq4_eq);
\end{matlabcode}
\begin{matlaboutput}
Soluzioni:
\end{matlaboutput}
\begin{matlabsymbolicoutput}
\hskip1em $\displaystyle \cos \left(q_4 \right)=\frac{A_{1,3} }{\sin \left(q_5 \right)}$
\end{matlabsymbolicoutput}
\begin{matlabsymbolicoutput}
\hskip1em $\displaystyle \sin \left(q_4 \right)=\frac{A_{2,3} }{\sin \left(q_5 \right)}$
\end{matlabsymbolicoutput}
\begin{matlabcode}
cosq6 = rhs(cosq6_eq);
sinq6 = rhs(sinq4_eq);
q4_val = atan2(sinq6, cosq6);
fprintf('Valore q4:');
disp(q4_val);
\end{matlabcode}
\begin{matlaboutput}
Valore q4:
\end{matlaboutput}
\begin{matlabsymbolicoutput}
\hskip1em $\displaystyle \textrm{atan2}\left(\frac{A_{2,3} }{\sin \left(q_5 \right)},\frac{A_{1,3} }{\sin \left(q_5 \right)}\right)$
\end{matlabsymbolicoutput}

\begin{flushleft}
    Si può definire un sistema con gli elementi $(3,1)$ e $(3,2)$ della matrice $^3 R_6$ per calcolare $q_6$
\end{flushleft}

\begin{matlabcode}
fprintf('Sistema Da Risolvere Per Trovare q6:');
eq_cosq6 = eqSystem(3,1);
eq_sinq6 = eqSystem(3,2);
disp(eq_cosq6);
disp(eq_sinq6);
\end{matlabcode}
\begin{matlaboutput}
Sistema Da Risolvere Per Trovare q6:
\end{matlaboutput}
\begin{matlabsymbolicoutput}
\hskip1em $\displaystyle -\cos \left(q_6 \right)\,\sin \left(q_5 \right)=A_{3,1} $
\end{matlabsymbolicoutput}
\begin{matlabsymbolicoutput}
\hskip1em $\displaystyle \sin \left(q_5 \right)\,\sin \left(q_6 \right)=A_{3,2} $
\end{matlabsymbolicoutput}
\begin{matlabcode}
fprintf('Soluzioni:');
cosq6_eq = isolate(eq_cosq6, cos(q6));
disp(cosq6_eq);
sinq6_eq = isolate(eq_sinq6, sin(q6));
disp(sinq6_eq);
\end{matlabcode}
\begin{matlaboutput}
Soluzioni:
\end{matlaboutput}
\begin{matlabsymbolicoutput}
\hskip1em $\displaystyle \cos \left(q_6 \right)=-\frac{A_{3,1} }{\sin \left(q_5 \right)}$
\end{matlabsymbolicoutput}
\begin{matlabsymbolicoutput}
\hskip1em $\displaystyle \sin \left(q_6 \right)=\frac{A_{3,2} }{\sin \left(q_5 \right)}$
\end{matlabsymbolicoutput}
\newpage
\begin{matlabcode}
cosq6 = rhs(cosq6_eq);
sinq6 = rhs(sinq6_eq);
q6_val = atan2(sinq6, cosq6);
fprintf('Valore q6:');
disp(q6_val);
\end{matlabcode}
\begin{matlaboutput}
Valore q6:
\end{matlaboutput}
\begin{matlabsymbolicoutput}
\hskip1em $\displaystyle \textrm{atan2}\left(\frac{A_{3,2} }{\sin \left(q_5 \right)},-\frac{A_{3,1} }{\sin \left(q_5 \right)}\right)$
\end{matlabsymbolicoutput}


\newpage
\matlabheading{Configurazioni Singolari}

\begin{par}
\begin{flushleft}
Una configurazione singolare si verifica quando la struttura del manipolatore limita la capacità di muoversi in certe direzioni, portando a una perdita di uno o più gradi di libertà . Dal punto di vista della cinematica inversa, le singolarità emergono quando le equazioni che usiamo per risolvere le variabili di giunto diventano indeterminate o non ammettono una soluzione unica.
\end{flushleft}
\end{par}

\vspace{1em}

\matlabheadingtwo{\textbf{Singolarità Della Struttura Portante}}

\begin{par}
\begin{flushleft}
Queste singolarità dipendono dai giunti $q_1 ,q_2 ,q_3$ e determinano la capacità del robot di posizionare il centro del polso. Si verifica quando:
\end{flushleft}
\end{par}

\matlabheadingthree{\textbf{Singolarità Di Bordo}}

\begin{par}
\begin{flushleft}
Si verifica quando l'argomento della radice quadrata usata per calcolare $q_3$ diventa nullo. Cioè quando
\end{flushleft}
\end{par}

\begin{par}
$\|\hat{d} {\|}^2 -d_2^2 =0$
\end{par}

\begin{par}
\begin{flushleft}
Dove $\|d{\|}^2$ è noto, poichè $\hat{d} =d_{06} -^0 R_6 \cdot d_1$ e i valori $d_{06}$ e $^0 R_6$ sono gli input del problema e $d_1$, $d_2$e $d_6$ sono i parametri fisici del sistema.
\end{flushleft}
\end{par}

\begin{par}
\begin{flushleft}
Se si verifica questo, il centro del polso si trova sul confine della superficie esterna del suo spazio di lavoro.
\end{flushleft}
\end{par}

\begin{matlabcode}
dHatNormSq_Sing = 0;

fprintf('Norma al quadrato del vettore calcolato (d_03):');
disp(norm_sq_d03);
\end{matlabcode}
\begin{matlaboutput}
Norma al quadrato del vettore calcolato (d_03):
\end{matlaboutput}
\begin{matlabsymbolicoutput}
\hskip1em $\displaystyle {d_2 }^2 +{q_3 }^2 $
\end{matlabsymbolicoutput}
\begin{matlabcode}
fprintf('Norma al quadrato del vettore target (d_hat): 0');
\end{matlabcode}
\begin{matlaboutput}
Norma al quadrato del vettore target (d_hat): 0
\end{matlaboutput}
\begin{matlabcode}
eq_q3_Sing = (dHatNormSq_Sing == norm_sq_d03);

fprintf('Equazione per trovare q3:');
disp(eq_q3_Sing);
\end{matlabcode}
\begin{matlaboutput}
Equazione per trovare q3:
\end{matlaboutput}
\begin{matlabsymbolicoutput}
\hskip1em $\displaystyle 0={d_2 }^2 +{q_3 }^2 $
\end{matlabsymbolicoutput}
\begin{matlabcode}

% Risolviamo per q3
sol_q3 = solve(eq_q3_Sing, q3, ReturnConditions=true);
fprintf('Valore Di q3:');
disp(sol_q3.q3);
\end{matlabcode}
\begin{matlaboutput}
Valore Di q3:
\end{matlaboutput}
\begin{matlabsymbolicoutput}
\hskip1em $\displaystyle 0$
\end{matlabsymbolicoutput}

\begin{par}
\begin{flushleft}
Cioè due soluzioni convergono in un'unica soluzione cioè $q_3 =0$. Si perde una delle due possibili configurazioni per raggiungere il punto target.
\end{flushleft}
\end{par}

\newpage
\matlabheadingthree{\textbf{Singolarità Di Allineamento}}

\begin{par}
\begin{flushleft}
Si verifica quando il centro del polso si trova sull'asse di rotazione del primo giunto (asse $Z$ della base). Matematicamente, questo significa che le componenti $(x,y)$ del vettore $d_{hat}$ di posizione desiderata del polso sono nulle, cioè
\end{flushleft}
\end{par}

\begin{par}
$d_{hat} =\left(\begin{array}{c}
0\\
0\\
d_z
\end{array}\right)$
\end{par}

\begin{par}
\begin{flushleft}
Poichè $d_{hat}$ è un vettore noto (è un input del problema) si può verificare direttamente se le sue prime due componenti sono nulle. Se il vettore è in questa configurazione, il braccio è orientato verticalmente puntando direttamente sopra o sotto la base del robot.
\end{flushleft}
\end{par}

\begin{par}
\begin{flushleft}
In questa configurazione, le equazioni per calcolare $q_1$ diventano indeterminate e quindi
\end{flushleft}
\end{par}

\begin{par}
$q_1 =atan2\bigl(\sin (q_1 ),\cos (q_1 )\bigr)=atan2\bigl(0,\;0\bigr)$
\end{par}

\begin{par}
\begin{flushleft}
Questo accade perché qualsiasi rotazione del giunto $q_1$ non altera la posizione del centro del polso, che rimane sull'asse $Z$. Esiste quindi un numero infinito di soluzioni per $q_1$, e il \textbf{manipolatore perde un grado di libertà } posizionale.
\end{flushleft}
\end{par}

\vspace{1em}
\matlabheadingtwo{\textbf{Singolarità Del Polso}}

\begin{par}
\begin{flushleft}
Questa singolarità dipende dai giunti $q_4 ,q_5 ,q_6$ e influenza la capacità del robot di orientare l'end-effector.
\end{flushleft}
\end{par}

\begin{par}
\begin{flushleft}
La singolarità si verifica quando $\sin (q_5 )=0$ che è equivalente a $\cos (q_5 )=\pm 1$. Poiché $\cos (q_5 )$ è determinato dall'elemento $(3,3)$ della matrice $^3 R_6$, la condizione da verificare è: $^3 R_6 (3,3)=\pm 1$.
\end{flushleft}
\end{par}

\begin{par}
\begin{flushleft}
Questa matrice è calcolata come
\end{flushleft}
\end{par}

\begin{par}
$^3 R_6 ={\bigl(^0 R_3 \bigr)}^{-1} \;^0 R_6 ={\bigl(^0 R_3 \bigr)}^T \;^0 R_6$
\end{par}

\begin{par}
\begin{flushleft}
Dove
\end{flushleft}
\end{par}

\begin{itemize}

   \item{\begin{flushleft} $^0 R_6$: è un dato di input del problema di cinematica inversa. \end{flushleft}}
   \item{\begin{flushleft} ${\bigl(^0 R_3 \bigr)}^T$ Dipende da $q_1 ,q_2 ,q_3$ che sono appena stati calcolati. \end{flushleft}}
\end{itemize}

\begin{par}
\begin{flushleft}
Nel caso di questa singolarità gli assi di rotazione del giunto $q_4$ e del giunto $q_6$ si allineano e quindi il polso perde un grado di libertà rotazionale e di conseguenza non può più ruotare indipendentemente attorno a tre assi ortogonali distinti.
\end{flushleft}
\end{par}

\begin{par}
\begin{flushleft}
Se si verifica questa singolarità è impossibile determinare in modo univoco i valori di $q_4$ e $q_6$ poichè equazioni per trovarli diventano dipendenti. Si può risolvere solo per la loro somma o differenza. Questo significa che esiste un numero infinito di coppie ($q_4$,$q_6$) che soddisfano l'orientamento target, e qualsiasi coppia è una soluzione valida.
\end{flushleft}
\end{par}

\begin{matlabcode}
R36_T_Sing = subs(R36_T, q5, 0);
fprintf('Matrice Rotazione R36 (Singolarita ):');
disp(simplify(R36_T_Sing));
\end{matlabcode}
\begin{matlaboutput}
Matrice Rotazione R36 (Singolarita ):
\end{matlaboutput}
\begin{matlabsymbolicoutput}
\hskip1em $\displaystyle \left(\begin{array}{ccc}
\cos \left(q_4 +q_6 \right) & -\sin \left(q_4 +q_6 \right) & 0\\
\sin \left(q_4 +q_6 \right) & \cos \left(q_4 +q_6 \right) & 0\\
0 & 0 & 1
\end{array}\right)$
\end{matlabsymbolicoutput}
\newpage
\begin{par}
\begin{flushleft}
\underline{Le variabili }$q_4$\underline{ e }$q_6$\underline{ non compaiono più singolarmente, ma solo all'interno della somma }$q_4 +q_6$\underline{. Questo dimostra che non si può separarle e calcolarle in modo univoco.}
\end{flushleft}
\end{par}

\begin{matlabcode}
syms theta real;
fprintf('Ponendo q4 + q6 = theta\n');
fprintf('Matrice Rotazione R36 (Singolarita ):');
R36_T_Sing = [cos(theta), -sin(theta) 0;
              sin(theta), cos(theta), 0;
              0         , 0         , 1]
\end{matlabcode}
\begin{matlaboutput}
Ponendo q4 + q6 = theta
Matrice Rotazione R36 (Singolarita ):
\end{matlaboutput}
\begin{matlabsymbolicoutput}
R36\_T\_Sing =

\hskip1em $\displaystyle \left(\begin{array}{ccc}
\cos \left(\theta \right) & -\sin \left(\theta \right) & 0\\
\sin \left(\theta \right) & \cos \left(\theta \right) & 0\\
0 & 0 & 1
\end{array}\right)$
\end{matlabsymbolicoutput}
\begin{matlabcode}
theta_val = atan2(R36_T_Sing(2,1), R36_T_Sing(1,1));
fprintf('Valore Di theta Calcolato:\n');
disp(theta_val);
\end{matlabcode}
\begin{matlaboutput}
Valore Di theta Calcolato:
\end{matlaboutput}
\begin{matlabsymbolicoutput}
\hskip1em $\displaystyle \textrm{atan2}\left(\sin \left(\theta \right),\cos \left(\theta \right)\right)$
\end{matlabsymbolicoutput}
\begin{matlabcode}
fprintf('Ponendo q4 = 0\n');
fprintf('theta = 0 + q6 \n');
fprintf('q6 = theta = \n');
disp(theta_val);
\end{matlabcode}
\begin{matlaboutput}
Ponendo q4 = 0
theta = 0 + q6
q6 = theta =
\end{matlaboutput}
\begin{matlabsymbolicoutput}
\hskip1em $\displaystyle \textrm{atan2}\left(\sin \left(\theta \right),\cos \left(\theta \right)\right)$
\end{matlabsymbolicoutput}

\newpage
\matlabheadingtwo{Calcoli Numerici}

\begin{matlabcode}
% Parametri Fisici Manipolatore
d2_num = 1;      % Lunghezza del braccio orizzontale
d6_num = 0.3;    % Lunghezza dell'utensile

% Posizione Desiderata End-Effector (INPUT)
d06_des = [0.8; 0.4; 0.7];

% Orientamento Desiderato End-Effector (INPUT)
theta_y = -pi/2;
R06_des = [cos(theta_y), 0, sin(theta_y);
           0,            1, 0;
           -sin(theta_y),0, cos(theta_y)];

fprintf('Parametri Robot: d2 = %.2f, d6 = %.2f\n', d2_num, d6_num);
fprintf('Posizione Desiderata (d06_des):\n');
disp(d06_des');
\end{matlabcode}
\begin{matlaboutput}
Parametri Robot: d2 = 1.00, d6 = 0.30
Posizione Desiderata (d06_des):
\end{matlaboutput}
\begin{matlaboutput}
    0.8000    0.4000    0.7000
\end{matlaboutput}
\begin{matlabcode}
fprintf('Orientamento Desiderato (R06_des):\n');
disp(R06_des);
\end{matlabcode}
\begin{matlaboutput}
Orientamento Desiderato (R06_des):
    0.0000         0   -1.0000
         0    1.0000         0
    1.0000         0    0.0000
\end{matlaboutput}
\begin{matlabcode}
%% Calcolo Cinematica Inversa Posizione (q1, q2, q3)

% 1. Calcoliamo Posizione Desiderata del Centro Polso (d_hat)
d1_num = [0; 0; d6_num];
d_hat_num = d06_des - R06_des * d1_num;
fprintf('Posizione Desiderata Centro Polso (d_hat):\n');
disp(d_hat_num');
\end{matlabcode}
\begin{matlaboutput}
Posizione Desiderata Centro Polso (d_hat):
\end{matlaboutput}
\begin{matlaboutput}
    1.1000    0.4000    0.7000
\end{matlaboutput}
\begin{matlabcode}

% 2. Calcolo Di q3
dHatNormSq_num = norm(d_hat_num)^2;

% Verifichiamo che la posizione sia raggiungibile
if dHatNormSq_num < d2_num^2
    error('Posizione Non Raggiungibile: Polso Troppo Vicino Alla Base.');
end

% Sostituiamo il valore numerico nell'equazione per q3
eq_q3_num = subs(eq_q3, [dHatNormSq, d2], [dHatNormSq_num, d2_num]);
sol_q3_num = solve(eq_q3_num, q3);

if numel(sol_q3_num) == 2
    % Caso standard: due soluzioni distinte (Gomito Alto e Basso)
    q3_pos_num = double(sol_q3_num(1));
    q3_neg_num = double(sol_q3_num(2));
    fprintf('Trovate due soluzioni per q3 (Gomito Alto/Basso): %.4f e %.4f\n', q3_pos_num, q3_neg_num);
\end{matlabcode}
\newpage
\begin{matlabcode}
elseif isscalar(sol_q3_num)
    % Caso speciale: una sola soluzione (singolarita di bordo)
    q3_pos_num = double(sol_q3_num);
    q3_neg_num = double(sol_q3_num);
    warning('Trovata una sola soluzione per q3 (configurazione singolare).');
else
    error('Nessuna soluzione reale trovata per q3.');
end
\end{matlabcode}
\begin{matlaboutput}
Trovate due soluzioni per q3 (Gomito Alto/Basso): -0.9274 e 0.9274
\end{matlaboutput}
\begin{matlabcode}
% Si Sceglie La Soluzione "Gomito Alto"
q3_sol = q3_pos_num;

% 3. Calcolo Di q2
cosq2_num = d_hat_num(3) / q3_sol;

% Si Sceglie La Soluzione Positiva
sinq2_num = sqrt(1 - cosq2_num^2);
q2_sol = atan2(sinq2_num, cosq2_num);

% 4. Calcolo Di q1
vars_to_sub = {d2, q3, q2, d_hat(1), d_hat(2)};
vals_to_sub = {d2_num, q3_sol, q2_sol, d_hat_num(1), d_hat_num(2)};

% sin(q1) = (d2*dx - dy*q3*sin(q2)) / (d2^2 + q3^2*sin(q2)^2)
% cos(q1) = (d2*dy + dx*q3*sin(q2)) / (d2^2 + q3^2*sin(q2)^2)
% dx, dy sono gli elementi di d_hat

num_sin = d2_num * d_hat_num(1) - d_hat_num(2) * q3_sol * sinq2_num;
num_cos = d2_num * d_hat_num(2) + d_hat_num(1) * q3_sol * sinq2_num;
den = d2_num^2 + (q3_sol * sinq2_num)^2;

sinq1_num = num_sin / den;
cosq1_num = num_cos / den;

q1_sol = atan2(sinq1_num, cosq1_num);

% Calcolo Cinematica Inversa Orientamento (q4, q5, q6)

% Calcolo Matrice R03
% La matrice R03 e la parte di rotazione di T03
R03_sol = [ cos(q1_sol)*cos(q2_sol), -sin(q1_sol),  cos(q1_sol)*sin(q2_sol);
            cos(q2_sol)*sin(q1_sol),  cos(q1_sol),  sin(q1_sol)*sin(q2_sol);
            -sin(q2_sol),             0,            cos(q2_sol)           ];

fprintf('Matrice R03');
disp(R03_sol);
\end{matlabcode}
\begin{matlaboutput}
Matrice R03
    0.1483   -0.9805   -0.1288
   -0.7401   -0.1964    0.6431
   -0.6559         0   -0.7548
\end{matlaboutput}
\newpage
\begin{matlabcode}
% Calcolo Matrice Rotazione Polso R36
% R36 = (R03)^T * R06_des
R36_sol = R03_sol' * R06_des;
fprintf('Matrice R36');
disp(R36_sol);
\end{matlabcode}
\begin{matlaboutput}
Matrice R36
   -0.6559   -0.7401   -0.1483
   -0.0000   -0.1964    0.9805
   -0.7548    0.6431    0.1288
\end{matlaboutput}
\begin{matlabcode}
% Calcolo Di q5
cosq5_num = R36_sol(3,3);

if abs(cosq5_num) > 1
    cosq5_num = sign(cosq5_num);
end

% Si Sceglie La Soluzione Positiva
sinq5_num = sqrt(1-cosq5_num^2);
q5_sol = atan2(sinq5_num, cosq5_num);

% Assumendo sin(q5) != 0
if abs(sinq5_num) < 1e-6
    warning('Singolarita Polso (sin(q5) = 0)!');
    % Si Pone Uno A Zero E Si Calcola L'Altro
    q4_sol = 0; % Scelta arbitraria
    cosq6_num = R36_sol(1,1);
    sinq6_num = -R36_sol(1,2);
    q6_sol = atan2(sinq6_num, cosq6_num);
else
    % Calcolo Di q4
    sinq4_num = R36_sol(2,3) / sinq5_num;
    cosq4_num = R36_sol(1,3) / sinq5_num;
    q4_sol = atan2(sinq4_num, cosq4_num);

    % Calcolo Di q6
    sinq6_num = R36_sol(3,2) / sinq5_num;
    cosq6_num = -R36_sol(3,1) / sinq5_num;
    q6_sol = atan2(sinq6_num, cosq6_num);
end

fprintf('SOLUZIONE CINEMATICA INVERSA:\n');

fprintf('q1 = %.3f rad = %4.2fdeg\n', q1_sol, rad2deg(q1_sol));
fprintf('q2 = %.3f rad = %4.2fdeg\n', q2_sol, rad2deg(q2_sol));
fprintf('q3 = %.3f m \n', q3_sol);
fprintf('q4 = %.3f rad = %4.2fdeg\n', q4_sol, rad2deg(q4_sol));
fprintf('q5 = %.3f rad = %4.2fdeg\n', q5_sol, rad2deg(q5_sol));
fprintf('q6 = %.3f rad = %4.2fdeg\n', q6_sol, rad2deg(q6_sol));
\end{matlabcode}
\begin{matlaboutput}
SOLUZIONE CINEMATICA INVERSA:
q1 = 1.769 rad = 101.33deg
q2 = 2.426 rad = 139.01deg
q3 = -0.927 m
q4 = 1.721 rad = 98.60deg
q5 = 1.442 rad = 82.60deg
q6 = 0.706 rad = 40.43deg
\end{matlaboutput}

\end{document}
